\begin{abstract}
Large language models increasingly operate in multi-step settings where strategic misrepresentation can improve local outcomes, yet most empirical work on honesty still centers on single-turn question answering or static preference probes. We introduce \textbf{Bullshit-Bench}, a reproducible benchmark release for studying deception and instruction compliance by having four LLM agents play the card game \emph{Bullshit}. The setting is useful because lying is a legal action, the truth of each claim is objectively verifiable from hidden cards, and outcomes depend on both bluffing and lie detection under uncertainty. We evaluate six hosted models under four prompt conditions: a low-strategy prompt control, a deception-allowed condition, an asymmetric-fairness condition, and a plain-language honesty mandate. The release combines a seeded TypeScript engine, a public reset/step environment API, baseline comparison policies, provenance-aware logging, long-run failure recovery, and a player-game analysis pipeline. The frozen comparable cohort contains 600 winner-terminated games with no mixed-cohort or turn-cap exclusions. The main empirical result is that explicit honesty instructions reduce deceptive play but do not eliminate it: mean model-level lie frequency falls from 39.6\% in the deception-allowed condition to 27.2\% under the honesty mandate, while mean challenge frequency falls from 45.3\% to 21.8\% and mean lie success rises from 11.0\% to 34.7\%. The asymmetric-fairness condition yields heterogeneous model-level lie-frequency shifts ranging from -7.4 to +3.3 percentage points rather than a universal restraint effect. These results show that prompt framing changes the enforcement regime of the table, not just the liar's behavior.
\end{abstract}
